\documentclass[tikz,border=0pt]{standalone}
\usepackage{tikz}
\usetikzlibrary{shadings,fadings,calc,decorations.pathmorphing}

\tikzfading[name=soft,    inner color=transparent!0,   outer color=transparent!100]
\tikzfading[name=fade,    inner color=transparent!100, outer color=transparent!0]
\tikzfading[name=vignette,inner color=transparent!100, outer color=transparent!0]

\begin{document}
	\begin{tikzpicture}[line cap=round, line join=round]
		\def\W{9}
		\def\H{16}

		\clip (0,0) rectangle (\W,\H);

		% ================= 1. 宣纸底色 =================
		\fill[white!97!yellow] (0,0) rectangle (\W,\H);

		% 纸纹：极淡的横向纤维
		\pgfmathsetseed{11}
		\foreach \i in {1,...,260} {%
			\pgfmathsetmacro{\px}{9*rnd}%
			\pgfmathsetmacro{\py}{16*rnd}%
			\pgfmathsetmacro{\pl}{0.3 + 1.2*rnd}%
			\pgfmathsetmacro{\op}{0.01 + 0.03*rnd}%
			\draw[gray!40!brown, opacity=\op, line width=0.25pt]
			(\px,\py) -- (\px+\pl,\py);
		}

		% 淡淡的做旧水渍
		\foreach \cx/\cy/\rx/\ry/\op in {%
			2.0/12.5/2.2/1.6/0.045,
			7.0/13.8/1.8/1.3/0.035,
			1.5/4.0/2.4/1.8/0.040,
			7.5/3.5/2.0/1.5/0.030,
			4.5/8.0/3.0/2.2/0.025%
		} {%
			\fill[brown!40!gray, path fading=soft, opacity=\op]
			(\cx,\cy) ellipse [x radius=\rx, y radius=\ry];
		}

		% ================= 2. 远山（淡墨，多座） =================
		% 最远一层，几乎看不见
		\fill[gray!55!blue, opacity=0.12]
		(0,9.6) --
		(0,9.8) -- (0.9,10.5) -- (1.8,10.0) -- (2.6,10.7) --
		(3.5,10.2) -- (4.5,10.9) -- (5.5,10.3) -- (6.4,10.8) --
		(7.2,10.1) -- (8.1,10.6) -- (9,10.0) --
		(9,9.6) -- cycle;

		% 远山第二层
		\fill[gray!60!blue, opacity=0.18]
		(0,9.0) --
		(0,9.3) -- (1.1,10.0) -- (2.0,9.5) -- (3.0,10.2) --
		(4.2,9.6) -- (5.2,10.3) -- (6.3,9.7) -- (7.3,10.2) --
		(8.2,9.6) -- (9,10.0) --
		(9,9.0) -- cycle;

		% 中山（主峰，偏左）
		\fill[gray!45!blue, opacity=0.30]
		(0,8.2) --
		(0,8.6) -- (0.7,9.2) -- (1.4,8.7) -- (2.2,9.8) --
		(3.0,9.0) -- (3.6,9.6) -- (4.2,9.1) --
		(5.0,8.4) -- (6.0,8.0) -- (7.0,7.6) -- (8.0,7.4) --
		(9,7.3) --
		(9,8.2) -- cycle;

		% 主峰上方的小尖
		\fill[gray!35!blue, opacity=0.35]
		(2.0,9.5) -- (2.2,9.8) -- (2.5,9.4) -- cycle;

		% 中山右边淡墨补充
		\fill[gray!65!blue, opacity=0.15]
		(4.5,8.5) --
		(4.5,8.8) -- (5.3,9.3) -- (6.1,8.9) -- (7.0,9.2) --
		(7.8,8.8) -- (9,9.1) --
		(9,8.5) -- cycle;

		% ================= 3. 水面（大面积留白 + 几笔横线） =================
		% 水岸线
		\draw[gray!55!black, opacity=0.28, line width=0.5pt]
		(0,7.3) .. controls (2.5,7.15) and (5.5,7.35) .. (9,7.2);

		% 水面横纹（近实远虚）
		\foreach \y/\xstart/\xend/\op/\lw in {%
			7.05/0.8/8.2/0.14/0.35,
			6.85/1.4/7.6/0.12/0.32,
			6.65/0.6/8.4/0.11/0.30,
			6.45/2.0/7.0/0.10/0.28,
			6.25/1.2/7.8/0.09/0.26,
			6.05/2.5/6.5/0.08/0.24,
			5.85/1.6/7.4/0.07/0.22,
			5.65/2.8/6.2/0.06/0.20,
			5.45/2.2/6.8/0.05/0.18%
		} {%
			\draw[gray!60!black, opacity=\op, line width=\lw pt]
			(\xstart,\y) -- (\xend,\y);
		}

		% ================= 4. 近处小舟 =================
		\begin{scope}[shift={(2.6,6.55)}, rotate=-3]
			% 船身
			\draw[black!75!gray, line width=0.7pt, opacity=0.75]
			(-0.45,0) .. controls (-0.25,-0.16) and (0.25,-0.16) .. (0.45,0);
			% 桅杆
			\draw[black!75!gray, line width=0.5pt, opacity=0.7]
			(0.0,0) -- (0.0,0.38);
			% 小帆
			\fill[black!70!gray, opacity=0.55]
			(0.0,0.36) -- (0.22,0.14) -- (0.0,0.14) -- cycle;
			% 人影
			\fill[black!80!gray, opacity=0.7]
			(-0.18,0.02) circle (0.035);
			\draw[black!80!gray, line width=0.5pt, opacity=0.65]
			(-0.18,-0.02) -- (-0.18,0.0);
		\end{scope}

		% 船的倒影
		\draw[black!70!gray, opacity=0.18, line width=0.5pt]
		(2.15,6.45) .. controls (2.6,6.38) and (3.05,6.42) .. (3.05,6.45);

		% ================= 5. 飞鸟 =================
		% 远处一群
		\foreach \bx/\by/\s/\op in {%
			5.8/11.6/1.0/0.45,
			6.4/11.9/0.85/0.40,
			6.9/11.5/0.75/0.38,
			6.1/12.2/0.7/0.35,
			7.4/11.8/0.6/0.30,
			5.3/12.0/0.55/0.28%
		} {%
			\draw[black!70!gray, line width=0.8pt*\s, opacity=\op]
			(\bx,\by)
			.. controls (\bx+0.16*\s,\by+0.10*\s) and (\bx+0.28*\s,\by+0.10*\s)
			.. (\bx+0.40*\s,\by)
			.. controls (\bx+0.52*\s,\by+0.10*\s) and (\bx+0.64*\s,\by+0.10*\s)
			.. (\bx+0.80*\s,\by);
		}

		% 近处两只，墨色更重
		\draw[black!80!gray, line width=1.4pt, opacity=0.70]
		(2.4,12.6)
		.. controls (2.65,12.85) and (2.85,12.85) .. (3.05,12.6)
		.. controls (3.25,12.85) and (3.45,12.85) .. (3.70,12.6);
		\draw[black!80!gray, line width=1.2pt, opacity=0.60]
		(3.3,13.1)
		.. controls (3.5,13.3) and (3.68,13.3) .. (3.85,13.1)
		.. controls (4.02,13.3) and (4.20,13.3) .. (4.40,13.1);

		% ================= 6. 近景浓墨竹枝（右下角） =================
		% 主枝干
		\draw[black!82!gray, line width=2.6pt, opacity=0.80]
		(9,4.2) .. controls (7.8,3.6) and (7.0,2.6) .. (6.6,1.2);
		\draw[black!82!gray, line width=1.8pt, opacity=0.75]
		(7.7,3.55) .. controls (7.0,3.3) and (6.5,3.1) .. (6.0,3.2);
		\draw[black!82!gray, line width=1.3pt, opacity=0.70]
		(7.2,2.85) .. controls (6.6,2.3) and (6.2,1.9) .. (5.8,2.0);
		\draw[black!82!gray, line width=0.9pt, opacity=0.60]
		(6.9,3.15) .. controls (6.5,3.6) and (6.3,3.9) .. (6.5,4.2);

		% 竹叶（细长笔触）
		\foreach \sx/\sy/\a/\l/\op in {%
			6.65/3.35/155/0.55/0.70,
			6.50/3.30/200/0.45/0.65,
			6.35/3.22/170/0.50/0.60,
			6.10/3.18/150/0.42/0.58,
			5.95/3.22/195/0.40/0.55,
			6.05/2.05/210/0.45/0.62,
			5.85/1.95/165/0.38/0.55,
			6.30/2.35/190/0.42/0.58,
			6.45/4.15/140/0.40/0.52,
			6.55/4.25/175/0.35/0.48%
		} {%
			\draw[black!78!gray, opacity=\op, line width=1.1pt,
			rotate around={\a:(6.6,2.8)}]
			(6.6,2.8) -- ++(\l,0);
		}

		% ================= 7. 朱红印章（左下角） =================
		\begin{scope}[shift={(1.35,2.1)}]
			% 印章外框（手绘感：略微不规则）
			\fill[red!75!black, opacity=0.88]
			(-0.42,-0.42) .. controls (-0.20,-0.46) and (0.22,-0.44) .. (0.42,-0.40)
			.. controls (0.46,-0.18) and (0.44,0.20) .. (0.40,0.42)
			.. controls (0.18,0.46) and (-0.20,0.44) .. (-0.42,0.40)
			.. controls (-0.46,0.18) and (-0.44,-0.20) .. (-0.42,-0.42)
			-- cycle;
			% 印章内白底
			\fill[white!97!yellow, opacity=0.95]
			(-0.30,-0.30) .. controls (-0.14,-0.33) and (0.16,-0.31) .. (0.30,-0.28)
			.. controls (0.33,-0.12) and (0.31,0.14) .. (0.28,0.30)
			.. controls (0.12,0.33) and (-0.14,0.31) .. (-0.30,0.28)
			.. controls (-0.33,0.12) and (-0.31,-0.14) .. (-0.30,-0.30)
			-- cycle;
			% 内白底上的红字（用简单笔画模拟）
			% 左竖
			\draw[red!75!black, line width=1.2pt, opacity=0.85]
			(-0.16,-0.16) -- (-0.16,0.18);
			% 横折
			\draw[red!75!black, line width=1.0pt, opacity=0.85]
			(-0.16,0.06) -- (0.10,0.06) -- (0.10,-0.10);
			% 右点
			\fill[red!75!black, opacity=0.85]
			(0.16,0.14) circle (0.035);
			% 下方小横
			\draw[red!75!black, line width=0.9pt, opacity=0.80]
			(-0.06,-0.20) -- (0.14,-0.20);
		\end{scope}

		% ================= 8. 暗角（很轻，模拟旧纸） =================
		\fill[brown!50!black, path fading=vignette, opacity=0.18]
		(0,0) rectangle (\W,\H);

	\end{tikzpicture}
\end{document}